\documentclass[a4paper]{article}
\usepackage{fontspec}
\usepackage{graphicx}
\usepackage{amsmath,amssymb}
\usepackage{hyperref}
\usepackage{minted}
\usepackage{multirow}
\usepackage{float}
\usepackage{todonotes}
\usepackage{forest}
\usepackage{enumitem}
\usepackage[ngerman]{babel}

\title{}
\date{}

\setmainfont{Latin Modern Roman}
\newfontfamily\ipafont{FreeSerif}
\newcommand{\ipa}[1]{{\ipafont[#1]}}

\begin{document}
    \maketitle
    
    
    \section{Deutsche Aussprache und IPA}
        
        Für \textbf{deutsche Aussprache} und \textbf{IPA-Zeichen} wie \ipa{fɛɐ̯ˈʔantvɔʁtlɪç} gibt es ein paar gute Lernwege.
        
        \subsection{1. Wörterbücher mit IPA}
            \begin{itemize}[leftmargin=1.5cm]
        \item \textbf{Wiktionary}\\
        Dort findet man bei sehr vielen deutschen Wörtern die IPA-Aussprache.\\
        Beispiel: \textit{verantwortlich} $\rightarrow$ \ipa{fɛɐ̯ˈʔantvɔʁtlɪç}
        
        \item \textbf{Duden Online}\\
        Ein sehr zuverlässiges deutsches Wörterbuch mit Ausspracheangaben und oft auch Audio.
            \end{itemize}
        
        \subsection{2. IPA selbst lernen}
            Das System heißt \textbf{International Phonetic Alphabet (IPA)}.\\
            Wenn du die Zeichen lernst, kannst du die Aussprache von vielen Wörtern genauer verstehen.
            
            Eine gute Methode ist:
            \begin{itemize}[leftmargin=1.5cm]
        \item eine IPA-Tabelle anschauen
        \item auf die Laute hören
        \item die Zeichen mit deutschen Beispielen verbinden
            \end{itemize}
        
        \subsection{3. YouTube und Aussprachevideos}
            Suche zum Beispiel nach:
            \begin{itemize}[leftmargin=1.5cm]
        \item German IPA pronunciation
        \item German phonetics
            \end{itemize}
            
            Dort werden oft solche Laute erklärt:
            \begin{itemize}[leftmargin=1.5cm]
        \item \ipa{ç} wie in \textit{ich}
        \item \ipa{x} wie in \textit{Bach}
        \item \ipa{ʁ} als deutsches \textit{R}
            \end{itemize}
        
        \subsection{4. Wichtiger Tipp}
            Für Deutschlernen ist es \textbf{nicht unbedingt nötig}, IPA perfekt zu beherrschen.\\
            Oft reicht schon:
            \begin{itemize}[leftmargin=1.5cm]
        \item Audio hören
        \item nachsprechen
        \item Wörter mehrfach wiederholen
            \end{itemize}
            
            IPA ist besonders nützlich, wenn du die Aussprache \textbf{sehr genau} lernen willst.
    
    
    \section{Ein paar wichtige IPA-Zeichen im Deutschen}
        
        \begin{tabular}{ll}
            \ipa{ç}   & wie in \textit{ich}                         \\
            \ipa{x}   & wie in \textit{Bach}                        \\
            \ipa{ʁ}   & deutsches \textit{R}                        \\
            \ipa{ʔ}   & Knacklaut, wie am Anfang von \textit{Apfel} \\
            \ipa{ɛ}   & wie in \textit{Bett}                        \\
            \ipa{iː}  & wie in \textit{Liebe}                       \\
            \ipa{uː}  & wie in \textit{gut}                         \\
            \ipa{ɔ}   & wie in \textit{offen}                       \\
            \ipa{aɪ̯} & wie in \textit{mein}                        \\
            \ipa{aʊ̯} & wie in \textit{Haus}
        \end{tabular}
    
    
    \section{Beispiel}
        \textbf{verantwortlich} \ipa{fɛɐ̯ˈʔantvɔʁtlɪç}\\
        Das Wort bedeutet auf Englisch: \textit{responsible}.
    
    % \bibliographystyle{plain}
    % \bibliography{/home/donkarlo/Dropbox/repo/shared_research/bibliography/bibliography.bib}
\end{document}